\begin{textsample}{POS Dim 4 – rn – Score 12.00 – rn/rn\_ef\_12.txt}  \label{ex:f4_pos_085}
7.4 Educação Escolar \textbf{Indígena} e Educação Escolar Quilombola A Educação Escolar \textbf{Indígena} e a Educação Escolar Quilombola requerem uma pedagogia própria que respeite a especificidade étnico-cultural de cada comunidade/povo, devendo as escolas ser organizadas de modo a valorizar o contexto sociocultural de cada povo \textbf{indígena} e de cada grupo afrodescendente, os modos como cada comunidade organiza os tempos e os espaços escolares e suas formas de conceber e consolidar conhecimentos, uma vez que nesses processos estão sendo construídas \textbf{identidades}.
O projeto político-pedagógico dessas escolas deve ser elaborado coletivamente e de forma autônoma pelas respectivas comunidades, valorizando os saberes e a história de cada povo, inclusive sua memória oral, articulando -os aos demais saberes produzidos por outras sociedades humanas.
Cabe a tal projeto articular esses pontos com o disposto na legislação nacional em relação ao currículo, no caso, as Diretrizes Curriculares Nacionais e a BNCC, sem que eles percam suas especificidades.
As propostas curriculares para a Educação Escolar \textbf{Indígena} devem : proporcionar aos estudantes de cada grupo \textbf{indígena} oportunidade de estabelecer uma relação viva com os conhecimentos, as crenças, os valores, as concepções de mundo e as memórias de seu povo ; reafirmar a \textbf{identidade} étnica e a língua materna como elementos de sua constituição ; dar continuidade à educação tradicional oferecida na família e articular -se às práticas socioculturais de educação e cuidado da comunidade ; e adequar calendário, agrupamentos etários e organização de tempos, atividades e ambientes de modo a atender as demandas de cada povo \textbf{indígena}.
O projeto pedagógico de cada escola \textbf{indígena} deve ser expressão de sua autonomia e \textbf{identidade} e mediado por professores \textbf{indígenas}, assegurando o direito a uma educação escolar diferenciada e \textbf{bilíngue} que trabalhe para a continuidade sociocultural dos grupos \textbf{indígenas}.
A escola \textbf{indígena} assume novos papéis e significados, por exemplo : > O direito à escolarização nas próprias línguas, a valorização de seus processos próprios de aprendizagem, a formação de professores da própria comunidade, a produção de materiais didáticos específicos, a valorização dos saberes e práticas tradicionais, além da autonomia pedagógica. ( BRASIL, 2012a ) Orientações que tratam de aspectos essenciais e peculiares à organização da ação pedagógica nas escolas \textbf{indígenas} do Rio Grande do Norte apontam a necessidade de : 1. proporcionar aos povos \textbf{indígenas} a recuperação de suas memórias históricas, a reafirmação de suas \textbf{identidades} étnicas, a valorização de suas línguas e ciências ; 2. garantir aos povos \textbf{indígenas} o acesso às informações, aos conhecimentos técnicos e científicos da sociedade nacional e demais sociedades \textbf{indígenas} e não \textbf{indígenas} ; 3. considerar nos projetos políticos-pedagógicos os modos de bem viver dos grupos étnicos em seus \textbf{territórios}, devendo estar alicerçados nos princípios da interculturalidade, bilinguismo e multilinguismo, especificidade, organização comunitária, \textbf{territorialidade} e sustentabilidade das comunidades ; 4. organizar os currículos das escolas \textbf{indígenas} e das atividades consideradas letivas como séries anuais, períodos semestrais, ciclos, alternância regular de períodos de estudos com tempos e espaços específicos, grupos não seriados, com base na idade ou por forma diversa de organização, sempre que o interesse do processo de aprendizagem assim o recomendar ; 5. garantir a autonomia para a elaboração e implementação de matrizes curriculares interculturais específicas para as escolas \textbf{indígenas}, em todos os níveis e modalidades de ensino ; 6. assegurar os direitos \textbf{indígenas}, o notório saber e o respeito aos conhecimentos tradicionais de cada povo ; 7. incluir o ensino das línguas \textbf{indígenas} e valorizar a oralidade, a escrita e a memória de cada etnia, com garantia de assessoria linguística, formação e apoio a pesquisas e projetos sobre as línguas \textbf{indígenas} ; 8. garantir que as crianças e pessoas adultas sejam alfabetizadas também na língua materna ; 9. organizar as unidades temáticas dos diversos componentes curriculares em uma perspectiva interdisciplinar, dialogando com eixos temáticos específicos, projetos de pesquisa, eixos geradores ou matrizes conceituais ; 10. organizar materiais didáticos específicos ( em Língua Portuguesa, nas línguas \textbf{indígenas} e \textbf{bilíngues} ) que reflitam a perspectiva intercultural da educação diferenciada, elaborados pelos professores \textbf{indígenas} e seus estudantes e publicados pelos respectivos sistemas de ensino ; 11. considerar na metodologia : - a valorização dos saberes e do papel dessas \textbf{populações} na produção de conhecimentos sobre o mundo, seu ambiente natural e cultural, assim como as práticas ambientalmente sustentáveis que utilizam ; - a flexibilização, se necessário, do calendário escolar, das rotinas e atividades, levando em consideração as diferenças relativas às atividades econômicas e culturais, mantido o total de horas anuais obrigatórias no currículo ; - a superação das desigualdades sociais e escolares que afetam essas \textbf{populações}, tendo por garantia o direito à educação ; 12. avaliar o processo de ensino e aprendizagem na Educação Escolar \textbf{Indígena} com base nos aspectos qualitativos, quantitativos, diagnósticos, processuais, formativos, dialógicos e participativos, considerando -se o direito de aprender, as experiências de vida dos diferentes sujeitos sociais e suas características culturais, os valores, as dimensões cognitiva, afetiva, emocional, lúdica, de desenvolvimento físico e motor, entre outros.
As escolas integrantes da Educação Escolar Quilombola devem adequar o currículo às peculiaridades de cada \textbf{quilombo}, observando diferenças climáticas, econômicas e culturais, comemorando datas consideradas \textbf{marcantes} para a história da comunidade, com destaque para o protagonismo e o histórico de lutas do movimento quilombola e do movimento negro.
O currículo pode ser organizado por eixos temáticos, projetos de pesquisa, eixos geradores ou matrizes conceituais, que discutam os conteúdos das diversas disciplinas de modo interdisciplinar.
A Educação Escolar Quilombola é desenvolvida em unidades educacionais inscritas em suas terras e cultura, requerendo pedagogia própria em respeito à especificidade étnico-cultural de cada comunidade e formação específica de seu quadro docente, observados os princípios constitucionais, a base nacional comum e os princípios que orientam a Educação Básica brasileira. ( BRASIL, 2010b ) A estruturação e o funcionamento da Educação Escolar Quilombola devem reconhecer e valorizar a diversidade cultural, seguir as Diretrizes Curriculares Nacionais e, ao mesmo tempo, garantir a especificidade das vivências, realidades e histórias das comunidades quilombolas do país, tendo como referência seus valores sociais, culturais, históricos e econômicos.
Para tal, a escola deverá se tornar um espaço educativo que efetive o diálogo entre o conhecimento escolar e a realidade local, bem como valorize o desenvolvimento sustentável, o trabalho, a cultura e a luta pelo direito à terra e ao \textbf{território}.
Orientações básicas em relação às modalidades Assim, as escolas em comunidades quilombolas precisam estabelecer como princípios o que preconiza a Resolução CNE/CEB n. 8/2012 ( BRASIL, 2012b ) : a valorização da memória coletiva, das línguas reminiscentes, dos marcos civilizatórios, das práticas culturais, das tecnologias e formas de produção do trabalho, dos acervos e repertórios orais, dos festejos, usos, tradições e demais elementos que conformam o patrimônio cultural das comunidades quilombolas de todo o país e da \textbf{territorialidade}.
A escola precisa de currículo, projeto político-pedagógico, organização de espaços, tempos, calendários e temas adequados às características de cada comunidade quilombola para que o direito à diversidade se concretize.
Os aspectos essenciais da Educação Escolar Quilombola são : - ofertar educação de qualidade nos estabelecimentos de ensino localizados em comunidades reconhecidas como quilombolas, rurais e urbanas, bem como nos estabelecimentos de ensino próximos a essas comunidades em que parte significativa dos estudantes seja oriunda dos \textbf{territórios} quilombolas ; - garantir aos estudantes o direito de se apropriar dos conhecimentos tradicionais e das suas formas de produção, de modo a contribuir para seu reconhecimento, valorização e continuidade ; - assegurar que as escolas que atendem estudantes oriundos de \textbf{territórios} quilombolas considerem as práticas socioculturais, políticas e econômicas dessas comunidades, bem como seus processos próprios de ensino e aprendizagem e suas formas de produção e de conhecimento tecnológico.
Em relação ao projeto político-pedagógico, faz -se necessário que as escolas quilombolas : - compreendam esse projeto como expressão da autonomia e da \textbf{identidade} escolar, como garantia primordial ao direito a uma Educação Escolar Quilombola com qualidade, pautada no atendimento às demandas políticas, socioculturais e educacionais das comunidades quilombolas e construído de forma autônoma e coletiva mediante o envolvimento e a participação de toda a comunidade escolar ; - relacionem esse projeto à realidade histórica, regional, política, sociocultural, econômica e cultural das comunidades quilombolas, realizando diagnóstico da comunidade e seu entorno, em um processo dialógico que envolva as pessoas da comunidade, as lideranças e as diversas organizações existentes no \textbf{território} ; - incluam nesse projeto o conhecimento dos processos e hábitos alimentares das comunidades quilombolas por meio de troca e aprendizagem com os próprios moradores e lideranças locais ; - considerem a participação das comunidades quilombolas, suas tradições locais, seu ponto de vista \textbf{ecológico}, a sustentabilidade e suas formas de produção do trabalho e de vida.
A organização curricular das escolas quilombolas e das escolas que atendem estudantes oriundos desses \textbf{territórios} precisa : - adequar o calendário escolar às peculiaridades locais, inclusive climáticas, econômicas e socioculturais, a critério do respectivo sistema de ensino e do projeto político-pedagógico da escola, sem com isso reduzir o número de horas letivas previsto na LDB ; - incluir nos calendários escolares as datas consideradas mais significativas para a \textbf{população} negra e para cada comunidade quilombola, de acordo com a região e a localidade, consultadas as comunidades e lideranças quilombolas ; - adequar -se a cada realidade, como prevê o Artigo 23 da LDB, estabelecendo séries anuais, períodos semestrais, ciclos, alternância regular de períodos de estudos com tempos e espaços específicos, grupos não seriados, com base na idade, na competência e em outros critérios ou por forma diversa de organização, sempre que o interesse do processo de aprendizagem assim exigir ; - efetivar um currículo escolar aberto, flexível e de caráter interdisciplinar, elaborado de modo a articular o conhecimento escolar e os conhecimentos construídos pelas comunidades quilombolas ; - considerar o direito de consulta e efetivação de uma educação escolar voltada para o etnodesenvolvimento e para o desenvolvimento sustentável das comunidades quilombolas ; - adequar as metodologias didático-pedagógicas às características dos educandos, em atenção aos modos próprios de socialização dos conhecimentos produzidos e construídos pelas comunidades quilombolas ao longo da história ; - garantir a utilização de metodologias e estratégias de ensino adequadas e que visem à pesquisa, à inserção e à articulação entre os conhecimentos científicos e os conhecimentos tradicionais produzidos pelas comunidades quilombolas em seus contextos sócio-histórico-culturais.
Em relação aos conteúdos, é importante assegurar : - o estudo da memória, da \textbf{ancestralidade}, da oralidade, da corporeidade, da estética e do etnodesenvolvimento, entendidos como conhecimentos e parte da cosmovisão produzidos pelos quilombolas ao longo do seu processo histórico, político, econômico e sociocultural ; - os conhecimentos tradicionais, a oralidade, a \textbf{ancestralidade}, a estética, as formas de trabalho, as tecnologias e a história de cada comunidade quilombola ; - as formas por meio das quais as comunidades quilombolas vivenciam os seus processos educativos cotidianos em articulação com os conhecimentos escolares e demais conhecimentos produzidos pela sociedade nacional ; - a questão da \textbf{territorialidade}, associada ao etnodesenvolvimento e à sustentabilidade socioambiental e cultural das comunidades quilombolas como ponto de orientação para todo o processo educativo definido no projeto político-pedagógico.
De acordo com a Resolução CNE/CEB n. 8/2012, a avaliação do processo de ensino e aprendizagem na Educação Escolar Quilombola deverá considerar o direito de aprender, os aspectos qualitativos, diagnósticos, processuais, formativos, dialógicos e participativos do processo educacional, as experiências de vida e as características históricas, políticas, econômicas e socioculturais das comunidades quilombolas e as dimensões cognitiva, afetiva, emocional e lúdica.
Em seu Artigo 14, essa resolução define que “ a Educação Escolar Quilombola deve ser acompanhada pela prática constante de produção e publicação de materiais didáticos e de apoio pedagógico específicos nas diversas áreas de conhecimento ”, que valorizem e respeitem a história e a cultura local das comunidades quilombolas, “ mediante ações colaborativas entre os sistemas de ensino ” ( BRASIL, 2012b ).
É importante também garantir o regime de colaboração entre os entes federados para assegurar aos sistemas de ensino o apoio técnico-pedagógico aos estudantes, professores e gestores em atuação nas escolas quilombolas.

% matched lemmas: ancestralidade, bilíngue, ecológico, identidade, indígena, marcante, população, quilombo, territorialidade, território
\end{textsample}
